\clearpage

\section{Template 변환}

\subsection{개요 및 설계}
Template 모듈은 정적 분석 파이프라인 내부에서 CPG(Code Property Graph)를 템플릿 분석에 적합한 형태로 변환하는 핵심 계층이다. 소스 코드는 일차적으로 CPG로 변환되는데, 이 CPG는 타입 정보, 제어 흐름, 데이터 흐름, 호출 관계 등 매우 풍부한 의미 정보를 포함한다. 그러나 CPG 원본 데이터는 표현이 장황하고(verbose), 그래프 구조가 복잡하며, 동일한 의미의 코드라도 표현 방식이 다를 수 있다는 한계가 있어 보안 패턴 매칭이나 템플릿 기반 분석에 직접 사용하기에는 부적합하다.

따라서 Template 모듈은 복잡한 CPG를 다음과 같이 가공하여 후속 단계(패턴 매칭, 취약점 탐지, 코드 생성 등)에서 단일하고 일관된 중간 표현(Intermediate Representation, IR)을 다룰 수 있도록 한다.

\begin{itemize}
    \item \textbf{정규화}: \texttt{<operator>.addition}과 같은 내부 이름을 \texttt{+}와 같은 표준 연산자로 변환하고, 배열 선언 등 불일치한 표현을 일관된 템플릿 노드 타입으로 통합한다.
    \item \textbf{AST 재구성}: 그래프 엣지로 표현된 관계 중 AST 엣지만을 추출하여 직관적인 계층적 트리 구조를 재구성한다.
    \item \textbf{타입 추론 및 보강}: 암시적인 타입 정보를 추론하고, 원본 코드 메타데이터를 부착한다.
    \item \textbf{평탄화(Flattening)}: 필요 시 트리를 그래프 알고리즘에 적합한 노드/엣지 리스트 형태로 재가공한다.
\end{itemize}

\paragraph{\textbf{설계 근거 및 전략}:}
\begin{itemize}
    \item \textbf{컴포넌트 분리}: 추출, 변환, 후처리, 평탄화 단계를 명확히 분리하여 유지보수성과 테스트 용이성을 확보하였다.
    \item \textbf{재귀 처리}: AST의 중첩 구조를 직관적으로 처리하기 위해 재귀적 순회 방식을 채택하였다.
    \item \textbf{Graceful Degradation}: 비정상적인 CPG 입력(타입 정보 부족, 필수 속성 누락)에 대해서도 \texttt{<unknown>} 처리나 자식 노드 승격 등을 통해 분석을 지속한다. 단, 구조적 결함에 대해서는 명시적 예외를 발생시킨다.
    \item \textbf{설정 기반 관리}: 연산자 매핑이나 표준 라이브러리 목록을 코드가 아닌 별도 설정 파일로 관리하여 확장성을 높였다.
\end{itemize}

\subsection{기능 및 컴포넌트 구성}
Template 모듈은 4단계의 파이프라인을 통해 CPG를 변환하며, 각 단계는 전담 컴포넌트가 수행한다.

\begin{enumerate}
    \item \textbf{TemplateExtractor (추출)}: CPG 그래프에서 AST 구조를 추출한다.
    \begin{itemize}
        \item CPG JSON-LD 구조 검증 및 파싱 (\texttt{@value} 래퍼 해제).
        \item AST 엣지 필터링 및 \texttt{nodeInfoMap}, \texttt{childrenMap} 구성.
        \item 루트 노드(자식으로 등장하지 않는 노드) 식별 및 \texttt{TreeNode} 재귀적 생성.
    \end{itemize}
    
    \item \textbf{TemplateConverter (변환)}: 추출된 트리를 표준 템플릿 노드(\texttt{TemplateNodes})로 변환한다.
    \begin{itemize}
        \item 노드 라벨(CALL, LOCAL, METHOD, CONTROL\_STRUCTURE) 기반 디스패치.
        \item 연산자 이름 정규화 및 제어 구조 표준화.
        \item \texttt{BinaryUnaryTypeWrapper}를 이용한 Bottom-up 타입 추론 (실패 시 \texttt{<unknown>}).
        \item 표준 라이브러리 호출과 사용자 정의 호출 구분.
    \end{itemize}
    
    \item \textbf{PostProcessor (후처리)}: 변환된 트리를 정제하고 보강한다.
    \begin{itemize}
        \item \texttt{nodeType}이 없는 유효하지 않은 노드 제거 및 자식 승계.
        \item CPG의 \texttt{CODE} 속성을 추출하여 각 노드에 부착.
        \item 배열 선언 및 할당 노드 병합, Translation Unit 분리.
    \end{itemize}
    
    \item \textbf{PlanationTool (평탄화)}: 트리를 평탄화된 그래프로 변환한다.
    \begin{itemize}
        \item 계층적 트리를 노드 리스트와 엣지 리스트(\texttt{from}, \texttt{to})로 변환.
        \item 결정적 결과(Determinism) 보장을 위한 ID 기반 정렬.
        \item 엣지 유효성 검증 (Dangling edge 탐지).
    \end{itemize}
\end{enumerate}

\subsection{실행 흐름}
Template 변환은 \texttt{buildTemplateArtifacts()} 함수를 진입점으로 하여 순차적으로 실행된다. Endpoint는 용도에 따라 트리 기반의 \texttt{templateResult} 또는 그래프 기반의 \texttt{flatten}을 선택하여 사용할 수 있다.

% (*@ [H] 옵션을 사용하여 그림을 현재 위치에 고정 @*)
\begin{figure}[H]
    \centering
    \begin{lstlisting}[language=TypeScript]
function buildTemplateArtifacts(root: CPGRoot) {
  const extractor = new TemplateExtractor();
  const converter = new TemplateConverter();
  const postProcessor = new PostProcessor();
  const planationTool = new PlanationTool([...]);

  const template = extractor.getTemplateTree(root.export);
  const converted = converter.convertTree(template);
  let templateResult = postProcessor.removeInvalidNodes(converted);
  templateResult = postProcessor.addCodeProperties(templateResult, root);
  const flatten = planationTool.flatten(templateResult);
  
  return { template, templateResult, textLines, flatten };
}
    \end{lstlisting}
    \caption{Template 변환 파이프라인 오케스트레이션}
    \label{fig:build_artifacts}
\end{figure}

\subsubsection{Step 1: 트리 추출 (TemplateExtractor)}
CPG 데이터에서 \texttt{@value} 래퍼를 해제하고 AST 엣지만을 필터링하여 재귀적으로 구문 트리를 구성한다. 데이터 흐름이나 제어 흐름 엣지를 배제하여 순환 구조를 방지한다.

\begin{figure}[H]
    \centering
    \begin{lstlisting}[language=TypeScript]
public getTemplateTree(cpg: unknown): TreeNode[] {
  if (!cpg || !("@value" in cpg)) return [];
  const inner = (cpg as ICPGRootExport)["@value"];
  const edges = inner.edges.filter(e => e.label === "AST");
  const nodes = inner.vertices;
  // Construct nodeInfoMap, identify root, create tree recursively
}
    \end{lstlisting}
    \caption{CPG 파싱 및 AST 추출 로직}
    \label{fig:template_extractor}
\end{figure}

\subsubsection{Step 2: 노드 변환 (TemplateConverter)}
추출된 \texttt{TreeNode}를 표준화된 \texttt{TemplateNodes}로 변환한다. 각 노드의 라벨에 따라 적절한 핸들러로 분기(Dispatch)하며, 오류 발생 시 노드 식별 정보를 포함한 예외를 발생시킨다.

\begin{figure}[H]
    \centering
    \begin{lstlisting}[language=TypeScript]
private dispatchConvert(node: TreeNode): TemplateNodes | undefined {
    switch (node.label) {
      case "CALL": return this.handleCall(node);
      case "CONTROL_STRUCTURE": return this.handleControlStructure(node);
      case "LOCAL": return this.handleLocal(node);
      case "METHOD": return this.handleMethod(node);
    default: return this.handleSkippedNodes(node);
  }
}
    \end{lstlisting}
    \caption{노드 변환 디스패처}
    \label{fig:dispatch_convert}
\end{figure}

\paragraph{연산자 정규화 (CALL)} \texttt{<operator>.*} 형태의 이름을 \texttt{BinaryExpressionOperatorMap} 등을 통해 표준 기호(\texttt{+}, \texttt{==})로 매핑하고 \texttt{BinaryExpression} 또는 \texttt{UnaryExpression}으로 변환한다.

\begin{figure}[H]
    \centering
    \begin{lstlisting}[language=TypeScript]
private handleCallOperators(node: TreeNode) {
  if (BinaryExpressionOperatorMap[node.name]) {
    return {
      nodeType: TemplateNodeTypes.BinaryExpression,
      operator: BinaryExpressionOperatorMap[node.name],
      type: BinaryUnaryTypeWrapper(node),
      children: this.convertedChildren(node.children)
    };
  }
  // Handle UnaryExpression similarly
}
    \end{lstlisting}
    \caption{연산자 노드 변환 로직}
    \label{fig:handle_call_operators}
\end{figure}

\paragraph{제어 구조 (Control Structure)} \texttt{CONTROL\_STRUCTURE\_TYPE} 속성을 확인하여 IF, WHILE 등으로 변환한다. IF 문의 경우 조건, True 블록, False 블록을 템플릿 포맷에 맞게 재배치하며, 자식 노드 수가 부족할 경우 예외를 발생시킨다.

\begin{figure}[H]
    \centering
    \begin{lstlisting}[language=TypeScript]
private handleControlStructure(node: TreeNode) {
  const type = node.properties.CONTROL_STRUCTURE_TYPE;
  switch (type) {
    case "IF":
      if (node.children.length < 2) throw new Error("Invalid IF structure");
      return { nodeType: TemplateNodeTypes.IfStatement, 
               children: restructureIfChildren(node.children) };
    // Handle WHILE, FOR, SWITCH similarly
  }
}
    \end{lstlisting}
    \caption{제어 구조 정규화 로직}
    \label{fig:handle_control_structure}
\end{figure}

\paragraph{변수 및 함수 처리} 변수 선언(\texttt{LOCAL})의 경우 \texttt{TYPE\_FULL\_NAME}을 분석하여 배열(\texttt{[]}), 포인터(\texttt{*})를 구분한다. 함수(\texttt{METHOD})는 파일 스코프 내의 항목만 처리하며, 본문 유무에 따라 선언(Declaration)과 정의(Definition)로 나눈다.

\begin{figure}[H]
    \centering
    \begin{lstlisting}[language=TypeScript]
// Variable declaration type analysis
if (typeFullName.includes("[")) return { nodeType: TemplateNodeTypes.ArrayDeclaration };
if (typeFullName.includes("*")) return { nodeType: TemplateNodeTypes.PointerDeclaration };
return { nodeType: TemplateNodeTypes.VariableDeclaration };
    \end{lstlisting}
    \caption{변수 선언 타입 분석 로직}
    \label{fig:variable_logic}
\end{figure}

\subsubsection{Step 3: 후처리 (PostProcessor)}
\texttt{removeInvalidNodes}를 통해 \texttt{nodeType}이 없는 노드를 제거(자식은 승계)하고,\\
\texttt{addCodeProperties}를 통해 CPG의 \texttt{CODE} 속성을 템플릿 노드에 부착하여 디버깅 및 리포팅을 지원한다.

\begin{figure}[H]
    \centering
    \begin{lstlisting}[language=TypeScript]
public addCodeProperties(nodes: TemplateNodes[], cpg: CPGRoot): TemplateNodes[] {
  return nodes.map(node => {
    const vertex = cpg.export["@value"].vertices.find(v => v.id["@value"] === node.id);
    return { ...node, code: vertex?.properties.CODE || "" };
  });
}
    \end{lstlisting}
    \caption{코드 속성 부착 로직}
    \label{fig:add_code_properties}
\end{figure}

\subsubsection{Step 4: 평탄화 (PlanationTool)}
트리를 순회하여 노드를 복제하고 부모-자식 관계를 \texttt{from}-\texttt{to} 엣지로 변환한다. 노드와 엣지를 ID 기준으로 정렬하여 결과의 결정성(determinism)을 보장하며, 존재하지 않는 노드를 가리키는 엣지가 없는지 검증한다.

\begin{figure}[H]
    \centering
    \begin{lstlisting}[language=TypeScript]
public flatten(astRoots: TemplateNodes[]): TemplateFlattenedGraph[] {
  return astRoots.map(root => {
    this.traverse(root);
    this.nodes.sort((a, b) => a.id - b.id);
    this.edges.sort((e1, e2) => e1.from + e1.to - (e2.from + e2.to));
    this.validateEdges();
    return { edges: this.edges.slice(), nodes: this.nodes.slice() };
  });
}
    \end{lstlisting}
    \caption{평탄화 및 그래프 생성 로직}
    \label{fig:flatten_logic}
\end{figure}

\subsection{설정}
Template 모듈의 확장성을 위해 연산자 매핑, 표준 라이브러리 목록 등은 코드가 아닌 별도의 설정 파일로 관리된다.

\begin{itemize}
    \item \textbf{BinaryExpression.ts}: CPG 이진 연산자를 표준 연산자로 매핑한다.
    \begin{figure}[H]
        \centering
        \begin{lstlisting}[language=TypeScript]
export const BinaryExpressionOperatorMap = {
  "<operator>.addition": "+",
  "<operator>.equals": "==",
  "<operator>.logicalAnd": "&&",
  // ... more operators
};
        \end{lstlisting}
        \caption{이진 연산자 매핑 설정}
        \label{fig:binary_map}
    \end{figure}
    
    \item \textbf{UnaryExpression.ts}: 단항 연산자 매핑을 정의한다.
    \begin{figure}[H]
        \centering
        \begin{lstlisting}[language=TypeScript]
export const UnaryExpressionOperatorMap = {
  "<operator>.postIncrement": "++",
  "<operator>.logicalNot": "!",
  "<operator>.indirection": "*",
  // ... more operators
};
        \end{lstlisting}
        \caption{단항 연산자 매핑 설정}
        \label{fig:unary_map}
    \end{figure}
    
    \item \textbf{Predefined.ts}: \texttt{stdin}, \texttt{NULL} 등의 식별자에 대해 고정 타입을 부여한다.
    \item \textbf{StandardLibCall.ts}: 표준 라이브러리 및 보안 관련 함수 레지스트리를 관리하여 사용자 정의 호출과 구분한다.
\end{itemize}